\chapter{Statistika Deskriptif}\label{ch:g10:stats}

Statistika merangkum kumpulan bilangan yang besar --- nilai, gaji,
suhu --- dengan segelintir penunjuk yang dipilih baik-baik. Bab ini
memperkenalkan dua keluarga penunjuk: yang menentukan \emph{pusat}
datanya (\omterm{def:g10:stats:mean}{rata-rata}, \omterm{def:g10:stats:median}{median}) dan yang mengukur \emph{sebarannya} (\omterm{def:g10:stats:quartiles}{jangkauan},
\omterm{def:g10:stats:quartiles}{kuartil}). Pembahasan yang lebih dalam, termasuk simpangan baku, menyusul pada
\cref{ch:g11:stat}.

\section{Data dan frekuensi relatif}

\begin{definition}[Data statistik]\label{def:g10:stats:series}
Sebuah \emph{data statistik}\index{data statistik} adalah daftar nilai
$x_1, x_2, \dots, x_N$ yang teramati pada populasi berukuran $N$. Bila sebuah
nilai $x$ muncul $n$ kali, $n$ disebut \emph{frekuensinya} dan hasil bagi
\[
f = \frac{n}{N}
\]
disebut \emph{frekuensi relatifnya}\index{frekuensi relatif} (sering
dinyatakan dalam persen). Frekuensi relatif semua nilainya berjumlah $1$.
\end{definition}

\begin{example}\label{ex:g10:stats:frequencies}
Nilai $20$ siswa pada kuis berskala $5$ butir:
\begin{center}
\renewcommand{\arraystretch}{1.2}
\begin{tabular}{c|cccccc}
nilai $x$ & $0$ & $1$ & $2$ & $3$ & $4$ & $5$ \\
\hline
frekuensi $n$ & $1$ & $2$ & $4$ & $6$ & $5$ & $2$ \\
\omterm{def:g10:stats:series}{frekuensi relatif} $f$ & $0.05$ & $0.10$ & $0.20$ & $0.30$ & $0.25$ & $0.10$ \\
\end{tabular}
\end{center}
Periksa: $1 + 2 + 4 + 6 + 5 + 2 = 20$ dan \omterm{def:g10:stats:series}{frekuensi relatifnya} berjumlah $1$.
\end{example}

\begin{omfigure}
\begin{tikzpicture}
\begin{axis}[
  width=8.4cm, height=5.2cm,
  ybar, bar width=14pt,
  xmin=-0.7, xmax=5.7, ymin=0, ymax=7.4,
  xtick={0,1,2,3,4,5},
  ytick={1,2,3,4,5,6},
  xlabel={nilai}, ylabel={frekuensi},
  tick label style={font=\small},
  label style={font=\small}]
  \addplot[fill=omDef!30, draw=omDef] coordinates
    {(0,1) (1,2) (2,4) (3,6) (4,5) (5,2)};
\end{axis}
\end{tikzpicture}

{\small Diagram batang nilai kuisnya: satu batang untuk setiap nilai, dengan
tinggi sama dengan frekuensinya.}
\end{omfigure}

\section{Ukuran pemusatan}

\begin{definition}[Rata-rata]\label{def:g10:stats:mean}
\emph{Rata-rata}\index{rata-rata} sebuah data adalah
\[
\bar x = \frac{x_1 + x_2 + \dots + x_N}{N}.
\]
Jika nilai $v_1, \dots, v_k$ muncul dengan frekuensi $n_1, \dots, n_k$,
bilangan yang sama dihitung sebagai \emph{rata-rata berbobot}
\[
\bar x = \frac{n_1 v_1 + n_2 v_2 + \dots + n_k v_k}{n_1 + n_2 + \dots + n_k}.
\]
\end{definition}

\begin{example}\label{ex:g10:stats:mean}
Untuk kuis pada \cref{ex:g10:stats:frequencies}:
\[
\bar x = \frac{1 \times 0 + 2 \times 1 + 4 \times 2 + 6 \times 3
+ 5 \times 4 + 2 \times 5}{20}
= \frac{0 + 2 + 8 + 18 + 20 + 10}{20}
= \frac{58}{20} = 2.9 .
\]
\end{example}

\begin{definition}[Median]\label{def:g10:stats:median}
\emph{Median}\index{median} adalah nilai yang membelah data yang sudah
\emph{terurut} menjadi dua paruh: paling sedikit separuh nilainya $\leq$
median, dan paling sedikit separuh nilainya $\geq$ median. Dalam praktiknya,
urutkan $N$ nilai itu;
\begin{itemize}
  \item jika $N$ ganjil, mediannya adalah nilai di tengah, pada urutan
        $\frac{N+1}{2}$;
  \item jika $N$ genap, ambil \omterm{prop:g10:coordgeom:midpoint}{titik tengah} dua nilai di tengah, yaitu pada
        urutan $\frac N2$ dan $\frac N2 + 1$.
\end{itemize}
\end{definition}

\begin{example}\label{ex:g10:stats:median}
Data terurut $7$ harga rumah (dalam ribuan): $120$, $150$, $160$,
$180$, $210$, $240$, $900$. \omterm{def:g10:stats:median}{Mediannya} adalah nilai ke-$4$, yaitu $180$.
\omterm{def:g10:stats:mean}{Rata-ratanya} $\frac{1960}{7} = 280$ --- lebih besar daripada $6$ di antara $7$
harga itu! Satu nilai ekstrem ($900$) menarik \omterm{def:g10:stats:mean}{rata-rata} jauh lebih kuat
daripada \omterm{def:g10:stats:median}{median}: \omterm{def:g10:stats:median}{median} bersifat \emph{kekar}, \omterm{def:g10:stats:mean}{rata-rata} tidak.
\end{example}

\section{Ukuran penyebaran}

\begin{definition}[Jangkauan, kuartil]\label{def:g10:stats:quartiles}
Untuk data yang sudah terurut:
\begin{itemize}
  \item \emph{jangkauan}\index{jangkauan} adalah selisih antara nilai
        terbesar dan nilai terkecil;
  \item \emph{kuartil pertama}\index{kuartil} $Q_1$ adalah nilai terkecil
        yang membuat paling sedikit seperempat ($25\%$) nilainya
        $\leq Q_1$; \emph{kuartil ketiga} $Q_3$ adalah nilai terkecil
        yang membuat paling sedikit tiga perempat ($75\%$) nilainya
        $\leq Q_3$;
  \item \emph{jangkauan antarkuartil}\index{jangkauan antarkuartil} $Q_3 - Q_1$ mengukur sebaran
        paruh tengah datanya.
\end{itemize}
\end{definition}

\begin{method}[Mencari kuartil]\label{met:g10:stats:quartiles}
Untuk data dengan $N$ nilai terurut:
\begin{enumerate}
  \item hitunglah $\frac N4$; bila hasilnya bukan \omterm{def:g10:numbers:sets}{bilangan bulat}, bulatkan
        \emph{ke atas}; $Q_1$ adalah nilai pada urutan itu;
  \item hitunglah $\frac{3N}{4}$; bulatkan ke atas bila perlu; $Q_3$ adalah
        nilai pada urutan itu.
\end{enumerate}
\end{method}

\begin{example}\label{ex:g10:stats:quartiles}
Ambil $11$ nilai terurut berikut
\[
2,\ 3,\ 5,\ 5,\ 6,\ 7,\ 8,\ 9,\ 9,\ 10,\ 12 .
\]
\omterm{def:g10:stats:median}{Median}: urutan $\frac{11+1}{2} = 6$, jadi \omterm{def:g10:stats:median}{mediannya} $7$.
$Q_1$: $\frac{11}{4} = 2.75$, dibulatkan ke atas menjadi $3$; nilai ke-$3$ adalah $5$.
$Q_3$: $\frac{33}{4} = 8.25$, dibulatkan ke atas menjadi $9$; nilai ke-$9$ adalah $9$.
\omterm{def:g10:stats:quartiles}{Jangkauannya}: $12 - 2 = 10$; \omterm{def:g10:stats:quartiles}{jangkauan antarkuartilnya}: $9 - 5 = 4$.
\end{example}

\begin{omfigure}
\begin{tikzpicture}[scale=0.85]
  \draw[-Stealth] (1,0) -- (13.4,0) node[below, font=\small] {nilai};
  \foreach \x in {2,4,6,8,10,12}
    { \draw (\x,-0.07) -- (\x,0.07);
      \node[below, font=\footnotesize] at (\x,-0.1) {\x}; }
  \draw[omDef, thick, fill=omDef!12] (5,0.5) rectangle (9,1.3);
  \draw[omDef, very thick] (7,0.5) -- (7,1.3);
  \draw[omDef, thick] (2,0.9) -- (5,0.9);
  \draw[omDef, thick] (9,0.9) -- (12,0.9);
  \draw[omDef, thick] (2,0.7) -- (2,1.1);
  \draw[omDef, thick] (12,0.7) -- (12,1.1);
  \node[above, font=\small] at (5,1.32) {$Q_1$};
  \node[above, font=\small] at (7,1.32) {median};
  \node[above, font=\small] at (9,1.32) {$Q_3$};
  \node[above, font=\small] at (2,1.12) {min};
  \node[above, font=\small] at (12,1.12) {max};
\end{tikzpicture}

{\small Diagram kotak \cref{ex:g10:stats:quartiles}: kotaknya membentang
sepanjang \omterm{def:g10:numbers:interval}{interval} antarkuartil $\intcc{Q_1}{Q_3}$ dan memuat paruh tengah
datanya; kedua sungutnya mencapai nilai ekstremnya.}
\end{omfigure}

\begin{example}[Membandingkan dua kelompok]\label{ex:g10:stats:comparison}
Dua kelas mengerjakan ujian yang sama (nilai dari $20$):
\begin{center}
\renewcommand{\arraystretch}{1.2}
\begin{tabular}{c|ccc}
 & \omterm{def:g10:stats:median}{median} & $Q_1$ & $Q_3$ \\
\hline
kelas A & $12$ & $10$ & $14$ \\
kelas B & $12$ & $6$ & $17$ \\
\end{tabular}
\end{center}
\omterm{def:g10:stats:median}{Mediannya} sama, sebarannya sangat berbeda: kelas A homogen (separuh
nilainya di $\intcc{10}{14}$), kelas B heterogen (paruh tengahnya
membentang sepanjang $\intcc{6}{17}$). Ukuran pemusatan saja tidak pernah
menceritakan seluruh kisahnya.
\end{example}

\begin{remark}
Baik \omterm{def:g10:stats:mean}{rata-rata} maupun \omterm{def:g10:stats:median}{median} tidak ada yang ``lebih baik'': \omterm{def:g10:stats:mean}{rata-rata} memakai
setiap nilai (dan diperlukan untuk menghitung jumlah keseluruhan), \omterm{def:g10:stats:median}{median}
bertahan terhadap nilai ekstrem. Ringkasan data yang sungguh-sungguh memberi
paling sedikit satu ukuran pemusatan \emph{dan} satu ukuran penyebaran.
\end{remark}

\section{Latihan}

\begin{exercise}[$\star$]\label{exo:g10:stats:1}
Berikut banyaknya buku yang dibaca dalam setahun oleh $15$ siswa:
\[
0,\ 1,\ 1,\ 2,\ 2,\ 2,\ 3,\ 3,\ 4,\ 4,\ 5,\ 6,\ 7,\ 8,\ 12 .
\]
Hitunglah \omterm{def:g10:stats:mean}{rata-rata}, \omterm{def:g10:stats:median}{median} dan \omterm{def:g10:stats:quartiles}{jangkauan} data ini.
\end{exercise}

\begin{exercise}[$\star$]\label{exo:g10:stats:2}
Sebuah dadu dilempar $50$ kali:
\begin{center}
\renewcommand{\arraystretch}{1.2}
\begin{tabular}{c|cccccc}
hasil & $1$ & $2$ & $3$ & $4$ & $5$ & $6$ \\
\hline
frekuensi & $7$ & $9$ & $8$ & $10$ & $6$ & $10$ \\
\end{tabular}
\end{center}
Hitunglah \omterm{def:g10:stats:series}{frekuensi relatif} setiap hasil dan \omterm{def:g10:stats:mean}{rata-rata} hasilnya.
\end{exercise}

\begin{exercise}[$\star$]\label{exo:g10:stats:3}
Carilah \omterm{def:g10:stats:median}{median}, $Q_1$ dan $Q_3$ dari data terurut
\[
1,\ 2,\ 4,\ 4,\ 5,\ 6,\ 6,\ 7,\ 8,\ 9,\ 9,\ 11 \qquad (N = 12),
\]
lalu gambarlah diagram kotaknya.
\end{exercise}

\begin{exercise}[$\star$]\label{exo:g10:stats:4}
\omterm{def:g10:stats:mean}{Rata-rata} $24$ nilai ujian adalah $11.5$. Siswa ke-$25$ mengikuti ujian itu
dan memperoleh $19$. Berapa \omterm{def:g10:stats:mean}{rata-rata} kelas yang baru?
\end{exercise}

\begin{exercise}[$\star\star$]\label{exo:g10:stats:5}
Seorang siswa mempunyai \omterm{def:g10:stats:mean}{rata-rata} $12$ setelah $5$ ujian yang berbobot
sama. Nilai berapa pada ujian ke-$6$ yang akan menaikkan \omterm{def:g10:stats:mean}{rata-ratanya} menjadi
$13$? Apakah \omterm{def:g10:stats:mean}{rata-rata} $14$ terjangkau (nilainya dari $20$)?
\end{exercise}

\begin{exercise}[$\star\star$]\label{exo:g10:stats:6}
Ciptakan sebuah data berisi $7$ nilai taknegatif yang \omterm{def:g10:stats:median}{mediannya} $10$ dan
\omterm{def:g10:stats:mean}{rata-ratanya} lebih besar daripada $20$; lalu satu lagi yang \omterm{def:g10:stats:median}{mediannya} $10$ dan
\omterm{def:g10:stats:mean}{rata-ratanya} lebih kecil daripada $6$. Mengapa \omterm{def:g10:stats:mean}{rata-rata} data semacam itu tidak
pernah lebih kecil daripada $\frac{40}{7}$? Apa yang diperlihatkan contoh ini?
\end{exercise}

\begin{exercise}[$\star\star$]\label{exo:g10:stats:7}
Di sebuah perusahaan, $9$ karyawannya masing-masing bergaji $1800$ per bulan
dan direkturnya bergaji $12000$.
\begin{enumerate}
  \item Hitunglah \omterm{def:g10:stats:mean}{rata-rata} dan \omterm{def:g10:stats:median}{median} gajinya.
  \item Ukuran mana yang paling baik memerikan gaji ``khas'' di sini? Mengapa?
\end{enumerate}
\end{exercise}

\begin{exercise}[$\star\star$]\label{exo:g10:stats:8}
Sebuah kelas dengan $12$ murid laki-laki mempunyai \omterm{def:g10:stats:mean}{rata-rata} tinggi $168$ cm;
$18$ murid perempuan di kelas yang sama mempunyai \omterm{def:g10:stats:mean}{rata-rata} tinggi $163$ cm.
Hitunglah \omterm{def:g10:stats:mean}{rata-rata} tinggi seluruh kelas. (Hati-hati: nilainya \emph{bukan}
\omterm{prop:g10:coordgeom:midpoint}{titik tengah} $168$ dan $163$ --- bobotilah \omterm{def:g10:stats:mean}{rata-ratanya} dengan ukuran kelompoknya.)
\end{exercise}

\begin{exercise}[$\star\star\star$]\label{exo:g10:stats:9}
Sebuah data berisi $N$ nilai $x_1, \dots, x_N$ mempunyai \omterm{def:g10:stats:mean}{rata-rata} $\bar x$.
Setiap nilainya diubah menjadi $y_i = a x_i + b$ untuk bilangan tetap $a$ dan $b$.
\begin{enumerate}
  \item Tunjukkan bahwa \omterm{def:g10:stats:mean}{rata-rata} data yang baru adalah $a \bar x + b$.
  \item Suhu selama sepekan, dalam derajat Celsius, mempunyai
        \omterm{def:g10:stats:mean}{rata-rata} $20$.
        Berapa \omterm{def:g10:stats:mean}{rata-ratanya} dalam derajat Fahrenheit
        ($F = 1.8\,C + 32$)?
\end{enumerate}
\end{exercise}

\section{Soal: Orang rata-rata itu tidak ada}

\begin{problem}[{Soal akhir pekan --- rata-rata melawan median: pencilan,
penyesuaian nilai, gudang di tepi jalan, dan kokpit yang dirancang
untuk siapa-siapa}]
\label{pb:g10:stats:1}
Pada tahun 1950 Angkatan Udara Amerika Serikat mengukur $4\,000$ penerbang
pada sepuluh ukuran tubuh lalu membangun kokpit bagi \emph{lelaki \omterm{def:g10:stats:mean}{rata-rata}}.
Kokpit itu ternyata pas --- seperti ditemukan seorang letnan lewat
penghitungan --- untuk hampir tidak seorang pun. Statistika merangkum
kerumunan dengan bilangan tunggal, dan setiap ringkasan mengatakan kebenaran
tentang sesuatu sekaligus berbohong tentang hal lain. Soal ini menyeret
\omterm{def:g10:stats:mean}{rata-rata} dan \omterm{def:g10:stats:median}{median} (\cref{def:g10:stats:mean}, \cref{def:g10:stats:median})
ke persidangan, membiarkan masing-masing memperlihatkan keunggulannya, lalu
menutup perkara kokpit itu.

\medskip
\noindent\textbf{Bagian I --- Dua ringkasan diadili.}
\begin{enumerate}
  \item Sebuah jalan menjual tujuh rumah, dalam ribuan euro:
        $200$, $210$, $220$, $230$, $240$, $250$ --- dan satu
        rumah besar seharga $1\,200$. Hitunglah \omterm{def:g10:stats:mean}{rata-rata} dan \omterm{def:g10:stats:median}{median}
        harganya. Bilangan mana yang memerikan ``sebuah rumah di jalan ini''?
  \item Buang rumah besar itu lalu hitung ulang keduanya. Sejauh mana
        masing-masing ringkasan bergeser? Nyatakan moralnya tentang pencilan.
  \item Sebuah tabel sensus: dari $100$ keluarga, $20$ tidak
        punya anak, $25$ punya satu, $30$ punya dua, $15$ punya
        tiga, dan $10$ punya empat. Hitunglah \omterm{def:g10:stats:mean}{rata-rata} banyaknya anak. Tidak ada
        keluarga yang ``punya $1.7$ anak'' --- lalu besaran nyata apa
        yang disandikan \omterm{def:g10:stats:mean}{rata-rata} itu? (Kalikan dengan $100$.)
  \item Untuk kelima belas nilai
        $4, 6, 7, 8, 9, 10, 10, 11, 12, 12, 13, 14, 15, 17, 19$:
        carilah \omterm{def:g10:stats:median}{median} dan \omterm{def:g10:stats:quartiles}{kuartil} $Q_1$, $Q_3$
        (\cref{met:g10:stats:quartiles}).
  \item Untuk data yang sama, berikan \omterm{def:g10:stats:quartiles}{jangkauan} dan
        \omterm{def:g10:stats:quartiles}{jangkauan antarkuartilnya}. Apa yang diukur masing-masing?
\end{enumerate}

\medskip
\noindent\textbf{Bagian II --- Penyesuaian nilai dan kelas gabungan.}
\begin{enumerate}[resume]
  \item Dari \cref{exo:g10:stats:9}: menggeser setiap nilai sebesar
        $+b$ menggeser \omterm{def:g10:stats:mean}{rata-rata}, \omterm{def:g10:stats:median}{median} dan \omterm{def:g10:stats:quartiles}{kuartil} sebesar $+b$; penskalaan
        dengan $a$ menskalakan semuanya dengan $a$. Apa yang terjadi pada
        \omterm{def:g10:stats:quartiles}{jangkauan} dan \omterm{def:g10:stats:quartiles}{jangkauan antarkuartil} di bawah pergeseran? Di bawah
        penskalaan? Berikan alasannya.
  \item Sebuah ujian berrata-rata $9$ (dari $20$) dan gurunya menginginkan
        \omterm{def:g10:stats:mean}{rata-rata} $12$. Dua penyesuaian diusulkan: menambah $3$ angka kepada
        semua orang, atau mengalikan setiap nilai dengan $\frac43$. Hitunglah
        nasib siswa bernilai $6$ dan siswa bernilai $18$ di bawah masing-masing
        penyesuaian. Penyesuaian mana yang menguntungkan siapa --- dan yang
        mana bermasalah dengan batas atasnya?
  \item Kelas A ($20$ siswa) berrata-rata $12$; kelas B ($30$
        siswa) berrata-rata $10$. Hitunglah \omterm{def:g10:stats:mean}{rata-rata} $50$
        siswa itu bersama-sama, lalu jelaskan mengapa nilainya bukan $11$
        (pelajaran \omterm{def:g10:stats:mean}{rata-rata} berbobot pada soal akhir pekan tentang \omterm{def:g10:stats:mean}{rata-rata}
        harmonik di jilid sebelumnya, kini di ruang kelas).
  \item Tunjukkan bahwa \omterm{def:g10:stats:median}{median} menolak permainan ini: hitunglah \omterm{def:g10:stats:median}{median}
        $A = \{1, 2, 3\}$ dan $B = \{0, 10\}$, lalu
        \omterm{def:g10:stats:median}{median} daftar gabungannya, lalu bandingkan dengan
        \omterm{def:g10:stats:mean}{rata-rata} kedua \omterm{def:g10:stats:median}{median} itu yang dibobot ukurannya.
  \item Seluncur indah: lima juri memberi nilai $5.2$, $5.3$, $5.4$,
        $5.5$, $9.9$. Hitunglah \omterm{def:g10:stats:mean}{rata-rata}, \omterm{def:g10:stats:median}{median}, dan
        \emph{\omterm{def:g10:stats:mean}{rata-rata} terpangkas} (buang yang terendah dan tertinggi, lalu
        rata-ratakan sisanya). Mengapa kebanyakan olahraga berjuri memangkas?
\end{enumerate}

\medskip
\noindent\textbf{Bagian III --- Membaca sebaran.}
\begin{enumerate}[resume]
  \item Dua pabrik sama-sama menghasilkan baut dengan \omterm{def:g10:stats:mean}{rata-rata} panjang
        $50$ mm. Pabrik A: $Q_1 = 49.9$, $Q_3 = 50.1$; pabrik
        B: $Q_1 = 48$, $Q_3 = 52$. Kamu memerlukan baut yang dapat saling
        dipertukarkan: pabrik mana, dan statistik mana yang memutuskan?
  \item Di sebuah klinik, \omterm{def:g10:stats:median}{median} waktu tunggunya $10$ menit tetapi
        \omterm{def:g10:stats:mean}{rata-rata} waktu tunggunya $25$. Bentuk data seperti apa yang
        dibocorkan pasangan ini? Susunlah kumpulan lima nilai dengan \omterm{def:g10:stats:median}{median}
        dan \omterm{def:g10:stats:mean}{rata-rata} persis seperti itu.
  \item ``Bayimu berada pada persentil ke-75 untuk tinggi badan'' ---
        apa arti kalimat itu? Dan mungkinkah setiap anak di sebuah kota
        berada ``di atas \omterm{def:g10:stats:mean}{rata-rata}''? Di atas \omterm{def:g10:stats:median}{median}? Jelaskan
        bedanya.
  \item Kelas A: min $4$, $Q_1\,8$, \omterm{def:g10:stats:median}{median} $11$, $Q_3\,13$, maks
        $17$. Kelas B: min $2$, $Q_1\,9$, \omterm{def:g10:stats:median}{median} $11$, $Q_3\,12$,
        maks $19$. Bandingkan: nilai yang khas, kehomogenan
        paruh tengahnya, dan keliaran nilai ekstremnya.
  \item Daftar periksa sang wartawan: sebuah tajuk berbunyi ``\omterm{def:g10:stats:mean}{rata-rata}
        gaji di TechCorp: $5\,000$ euro''. Daftarkan ketiga
        pertanyaan yang diajarkan soal ini untuk ditanyakan sebelum
        mempercayai bahwa tajuk itu mengatakan sesuatu tentang karyawan
        yang khas.
\end{enumerate}

\medskip
\noindent\textbf{Bagian IV --- Perkara kokpit.}
\begin{enumerate}[resume]
  \item Bangunlah buktimu sendiri: susunlah lima harga rumah
        dengan \omterm{def:g10:stats:mean}{rata-rata} $300$ dan \omterm{def:g10:stats:median}{median} $200$ (dalam ribuan euro).
  \item \omterm{def:g10:stats:mean}{Rata-rata} adalah titik seimbangnya: buktikan dari definisinya
        bahwa simpangan terhadap \omterm{def:g10:stats:mean}{rata-rata} selalu saling menghapus,
        $\sum (x_i - \bar x) = 0$, lalu periksalah pada kumpulan datamu
        di pertanyaan 16.
  \item \omterm{def:g10:stats:median}{Median} adalah peminimum jarak tempuh: toko-toko berada pada
        kilometer $1$, $3$, $4$, $9$, $13$ sebuah jalan, dan sebuah
        gudang harus meminimumkan \emph{jarak total} ke kelima
        toko itu. Hitunglah jarak total itu untuk gudang di
        \omterm{def:g10:stats:median}{median} ($4$), di \omterm{def:g10:stats:mean}{rata-rata} ($6$), dan di $5$. Siapa yang menang?
        (Hal ini berlaku umum: \omterm{def:g10:stats:median}{median} meminimumkan jarak mutlak total ---
        ujilah tempat mana pun yang lain.)
  \item Perkara kokpit, terpecahkan: andaikan seorang penerbang mempunyai
        peluang $30\,\%$ untuk ``dekat dengan \omterm{def:g10:stats:mean}{rata-rata}'' pada satu ukuran
        tubuh, saling bebas untuk $10$ ukuran.
        Taksirlah peluang untuk dekat dengan \omterm{def:g10:stats:mean}{rata-rata} pada
        \emph{kesepuluh ukuran} --- lalu cocokkan hasilnya dengan apa yang
        ditemukan sang letnan pada $4\,000$ penerbang. Apa yang dibangun
        Angkatan Udara sebagai ganti kokpit \omterm{def:g10:stats:mean}{rata-rata} itu?
  \item Penutup --- panduan pemakainya, satu baris untuk masing-masing: bila
        yang penting adalah jumlah keseluruhan (anggaran, hasil panen), pakailah
        \dots; bila yang penting adalah individu yang khas (gaji, waktu
        tunggu), pakailah \dots; bila yang penting adalah keadilan sebarannya,
        laporkanlah \dots; bila pencemaran data mengancam, pakailah
        \dots. Tutuplah perkaranya dengan judul soal ini.

\end{enumerate}
\end{problem}
